\documentclass[12pt,a4paper]{article}

% ── Packages ────────────────────────────────────────────────────────────────
\usepackage[utf8]{inputenc}
\usepackage[T1]{fontenc}
\usepackage[english]{babel}
\usepackage{geometry}
\geometry{margin=2.5cm}

\usepackage{graphicx}
\usepackage{float}
\usepackage{caption}
\usepackage{subcaption}
\usepackage{hyperref}
\usepackage{listings}
\usepackage{xcolor}
\usepackage{booktabs}
\usepackage{tabularx}
\usepackage{fancyhdr}
\usepackage{parskip}
\usepackage{lmodern}
\usepackage{array}

% ── Page style ───────────────────────────────────────────────────────────────
\pagestyle{fancy}
\fancyhf{}
\rhead{TCP Camera System -- Unity \& Python}
\lhead{\leftmark}
\rfoot{\thepage}

% ── Code listing style ───────────────────────────────────────────────────────
\definecolor{codebg}{RGB}{245,245,245}
\definecolor{codeframe}{RGB}{200,200,200}
\definecolor{keyword}{RGB}{0,0,200}
\definecolor{comment}{RGB}{80,120,80}
\definecolor{string}{RGB}{160,32,32}

\lstdefinestyle{csharp}{
    language=[Sharp]C,
    backgroundcolor=\color{codebg},
    frame=single,
    rulecolor=\color{codeframe},
    basicstyle=\ttfamily\footnotesize,
    keywordstyle=\color{keyword}\bfseries,
    commentstyle=\color{comment}\itshape,
    stringstyle=\color{string},
    numbers=left,
    numberstyle=\tiny\color{gray},
    numbersep=8pt,
    tabsize=4,
    breaklines=true,
    showstringspaces=false,
    captionpos=b
}

\lstdefinestyle{python}{
    language=Python,
    backgroundcolor=\color{codebg},
    frame=single,
    rulecolor=\color{codeframe},
    basicstyle=\ttfamily\footnotesize,
    keywordstyle=\color{keyword}\bfseries,
    commentstyle=\color{comment}\itshape,
    stringstyle=\color{string},
    numbers=left,
    numberstyle=\tiny\color{gray},
    numbersep=8pt,
    tabsize=4,
    breaklines=true,
    showstringspaces=false,
    captionpos=b
}

% ── Color boxes for UDP/TCP comparison ───────────────────────────────────────
\definecolor{udpcolor}{RGB}{220,240,255}
\definecolor{tcpcolor}{RGB}{220,255,220}

% ── Title ────────────────────────────────────────────────────────────────────
\title{
    \vspace{-1cm}
    \Large\textbf{TCP-Based JPEG Capture in Unity for Remote Control}\\[0.5em]
    \large Laboratory Report -- TCP Version
}
\author{Miguel Rodriguez}
\date{April 2, 2026}

% ═══════════════════════════════════════════════════════════════════════════════
\begin{document}
\maketitle
\thispagestyle{fancy}

% ── Abstract ──────────────────────────────────────────────────────────────────
\begin{abstract}
This report presents a TCP-based evolution of the UDP JPEG streaming system developed
in a previous laboratory session. The system captures JPEG frames from a Unity camera
and streams them to a Python client over a persistent TCP connection, supporting
high-resolution images (1280$\times$720), WASD-based remote camera control, and
multi-client operation over a Local Area Network (LAN). A comparative analysis of UDP
and TCP is included, discussing the trade-offs between simplicity, reliability,
and latency for real-time image streaming applications.
\end{abstract}

\tableofcontents
\newpage

% ═══════════════════════════════════════════════════════════════════════════════
\section{Project Objective}

This project extends the UDP snapshot system by replacing the transport layer with TCP,
addressing the key limitations of UDP for image streaming:

\begin{itemize}
    \item \textbf{No datagram size limit:} UDP is restricted to $\approx$65\,KB per
          datagram. TCP is stream-based and imposes no message-size limit, enabling
          high-resolution JPEG frames (1280$\times$720 at quality 70 produces frames
          well within practical limits but scalable beyond 65\,KB).
    \item \textbf{Guaranteed delivery:} TCP ensures every byte arrives in order,
          eliminating silent frame loss present in UDP on real networks.
    \item \textbf{Persistent connection:} the client connects once and maintains the
          session, reducing per-request overhead.
\end{itemize}

The request string format remains identical to the UDP version:
\begin{center}
\texttt{w=1280;h=720;q=70;fps=10;move=S}
\end{center}

% ═══════════════════════════════════════════════════════════════════════════════
\section{UDP vs.\ TCP -- Comparative Analysis}

\subsection{Protocol Overview}

\begin{table}[H]
\centering
\renewcommand{\arraystretch}{1.4}
\begin{tabularx}{\textwidth}{>{\bfseries}p{3.5cm} X X}
\toprule
\textbf{Property} & \textbf{UDP} & \textbf{TCP} \\
\midrule
Connection model  & Connectionless (stateless) & Connection-oriented (stateful) \\
Delivery guarantee & None -- packets may be lost & Guaranteed -- retransmission on loss \\
Ordering          & Not guaranteed & Guaranteed in-order delivery \\
Message framing   & Built-in (datagrams) & Manual -- requires length prefix \\
Max payload       & $\approx$65\,KB per datagram & Unlimited (stream) \\
Latency           & Lower (no handshake, no ACK) & Higher (handshake + ACK overhead) \\
Fragmentation risk & Yes -- lost fragment = lost frame & None -- OS handles reassembly \\
Multi-client      & All share one socket & One socket per client \\
Reconnection      & Automatic (stateless) & Must re-establish TCP handshake \\
Implementation complexity & Low & Medium \\
\bottomrule
\end{tabularx}
\caption{Side-by-side comparison of UDP and TCP properties.}
\end{table}

\subsection{Strengths and Weaknesses}

\begin{table}[H]
\centering
\renewcommand{\arraystretch}{1.4}
\begin{tabularx}{\textwidth}{>{\bfseries}p{2cm} X X}
\toprule
 & \textbf{UDP} & \textbf{TCP} \\
\midrule
\textbf{Strengths}
& \vspace{-0.4cm}\begin{itemize}
    \item Minimal latency
    \item Simple implementation
    \item No connection overhead
    \item Ideal for localhost tests
  \end{itemize}
& \vspace{-0.4cm}\begin{itemize}
    \item No frame loss
    \item No size limit (HD+ images)
    \item Supports multiple clients
    \item Auto-reconnect in client
  \end{itemize} \\
\textbf{Weaknesses}
& \vspace{-0.4cm}\begin{itemize}
    \item 65\,KB payload limit
    \item Silent frame loss on LAN
    \item IP fragmentation risk
    \item No delivery confirmation
  \end{itemize}
& \vspace{-0.4cm}\begin{itemize}
    \item Higher latency (ACK/retransmit)
    \item Requires message framing
    \item Connection management needed
    \item Head-of-line blocking
  \end{itemize} \\
\textbf{Best use}
& Low-latency, loss-tolerant streaming (e.g.\ game telemetry, sensor feeds)
& Reliable image delivery, high-resolution frames, production LAN systems \\
\bottomrule
\end{tabularx}
\caption{Strengths and weaknesses of UDP vs.\ TCP for JPEG streaming.}
\end{table}

\subsection{Why TCP for This Project}

Given that the system transmits full JPEG frames at 1280$\times$720 resolution, TCP is
the more appropriate choice: a single lost IP fragment in UDP silently discards the
entire frame with no indication to the application. TCP's retransmission mechanism
prevents this at the cost of slightly higher latency, which is acceptable for a
10\,FPS control stream.

% ═══════════════════════════════════════════════════════════════════════════════
\section{System Architecture}

\subsection{Message Framing Protocol}

Unlike UDP (where each \texttt{Send()} call produces exactly one datagram), TCP is a
byte stream with no built-in message boundaries. A 4-byte big-endian length prefix is
prepended to every message in both directions:

\begin{center}
\texttt{[ 4 bytes: length (big-endian) ][ N bytes: payload ]}
\end{center}

\begin{description}
    \item[Request (client $\to$ server)] UTF-8 string, e.g.\ \texttt{w=1280;h=720;q=70;fps=10;move=S}
    \item[Response (server $\to$ client)] Raw JPEG bytes. Length\,=\,0 means no frame
    available yet.
\end{description}

\subsection{Communication Flow}

\begin{figure}[H]
\centering
\begin{verbatim}
  Python Client (Ubuntu)              Unity Server (Windows 11)
  ──────────────────────              ─────────────────────────────────────
  sock.connect(10.50.33.6:5000)  →   TcpListener.AcceptTcpClientAsync()
                                          spawns HandleClient(Task)

  Loop:
    send_framed(request_string)   →   ReadExactAsync(4 bytes) → parse length
                                      ReadExactAsync(N bytes) → parse params
                                      ParseAndStoreParams()
                                        lock(_paramLock) → FrameProducerLoop
                                        lock(_moveLock)  → Update()
    recv_framed()                 ←   WriteAsync(4 bytes length + JPEG)
    np.frombuffer + cv2.imdecode
    cv2.imshow()
\end{verbatim}
\caption{TCP communication flow between client and server.}
\end{figure}

% ═══════════════════════════════════════════════════════════════════════════════
\section{Unity Server -- TcpJpegSnapshotServer.cs}

\subsection{TCP Listener and Multi-Client Support}

\begin{lstlisting}[style=csharp, caption={Starting the TCP listener bound to all interfaces.}]
private void StartTcpServer()
{
    _cts      = new CancellationTokenSource();
    _listener = new TcpListener(IPAddress.Any, listenPort);
    _listener.Start();
    _ = Task.Run(() => AcceptLoop(_cts.Token));
}

private async Task AcceptLoop(CancellationToken token)
{
    while (!token.IsCancellationRequested)
    {
        TcpClient client = await _listener.AcceptTcpClientAsync();
        // Each client runs in its own Task -- supports simultaneous connections
        _ = Task.Run(() => HandleClient(client, token));
    }
}
\end{lstlisting}

\subsection{Framed Read/Write}

\begin{lstlisting}[style=csharp, caption={Reading a length-prefixed message from the TCP stream.}]
// Read exactly 'count' bytes -- TCP may deliver them in multiple chunks
private async Task<int> ReadExactAsync(NetworkStream stream,
                                        byte[] buf, int count,
                                        CancellationToken token)
{
    int total = 0;
    while (total < count)
    {
        int n = await stream.ReadAsync(buf, total, count - total, token);
        if (n == 0) break; // connection closed
        total += n;
    }
    return total;
}
\end{lstlisting}

\begin{lstlisting}[style=csharp, caption={Sending the JPEG response with length prefix.}]
byte[] lenOut = BitConverter.GetBytes(
                    IPAddress.HostToNetworkOrder(jpg.Length));
await stream.WriteAsync(lenOut, 0, 4, token);   // 4-byte header
await stream.WriteAsync(jpg,    0, jpg.Length, token); // JPEG payload
\end{lstlisting}

\subsection{WASD Movement}

The movement command is consumed and reset every \texttt{Update()} frame, so the
object only moves while the client actively sends a direction key:

\begin{lstlisting}[style=csharp, caption={Movement applied and reset on the Unity main thread.}]
private void Update()
{
    string move;
    lock (_moveLock) { move = _currentMove; _currentMove = "0"; }
    Vector3 dir = Vector3.zero;
    switch (move)
    {
        case "W": dir =  transform.forward; break;
        case "S": dir = -transform.forward; break;
        case "A": dir = -transform.right;   break;
        case "D": dir =  transform.right;   break;
    }
    transform.position += dir * moveSpeed * Time.deltaTime;
}
\end{lstlisting}

% ═══════════════════════════════════════════════════════════════════════════════
\section{Python Client -- tcp\_client\_viewer.py}

\subsection{Framing Helpers}

\begin{lstlisting}[style=python, caption={Send and receive framed messages over TCP.}]
def send_framed(sock, data: bytes):
    header = struct.pack(">I", len(data))   # 4-byte big-endian length
    sock.sendall(header + data)

def recv_exact(sock, n: int) -> bytes:
    buf = b""
    while len(buf) < n:
        chunk = sock.recv(n - len(buf))
        if not chunk:
            raise ConnectionError("Connection closed by server.")
        buf += chunk
    return buf

def recv_framed(sock) -> bytes:
    length = struct.unpack(">I", recv_exact(sock, 4))[0]
    return recv_exact(sock, length) if length > 0 else b""
\end{lstlisting}

\subsection{Auto-Reconnect}

Unlike the UDP client, the TCP client detects disconnection and automatically
reconnects without restarting the script:

\begin{lstlisting}[style=python, caption={Outer reconnect loop wrapping the streaming loop.}]
while True:
    sock = socket.socket(socket.AF_INET, socket.SOCK_STREAM)
    try:
        sock.connect((args.ip, args.port))
    except (ConnectionRefusedError, socket.timeout):
        time.sleep(2)
        continue

    try:
        while True:          # streaming loop
            send_framed(sock, payload)
            data = recv_framed(sock)
            # decode + display ...
    except (ConnectionError, OSError):
        time.sleep(2)        # reconnect on any network error
    finally:
        sock.close()
\end{lstlisting}

% ═══════════════════════════════════════════════════════════════════════════════
\section{Simulation Results}

The system was tested in a two-machine LAN scenario:
\textbf{Unity server} on Windows\,11 at \texttt{10.50.33.6},
\textbf{Python client} on Ubuntu (Lenovo ThinkPad) at \texttt{10.50.33.4}.
Resolution was set to 1280$\times$720 with JPEG quality 70 at 10\,FPS.

\subsection{LAN Setup}

\begin{figure}[H]
    \centering
    \includegraphics[width=0.92\textwidth]{lab_setup_tcp.png}
    \caption{Physical laboratory setup. \textbf{Left:} Lenovo ThinkPad running Ubuntu
             (Python TCP client) showing the snapshot window and terminal.
             \textbf{Right:} Desktop PC running Windows\,11 with Unity\,6 in Play mode
             (TCP server) displaying the Game view and Console logs.
             Both machines share the \texttt{central.local} Wi-Fi network.}
    \label{fig:lab_setup}
\end{figure}

\subsection{Server -- Console and Network Configuration}

\begin{figure}[H]
    \centering
    \includegraphics[width=0.92\textwidth]{server_console_tcp.png}
    \caption{Unity Console and PowerShell \texttt{ipconfig} running simultaneously.
             The Console logs show repeated TCP requests from \texttt{10.50.33.4}
             (Ubuntu client) with the full parametric string
             \texttt{w=1280;h=720;q=70;fps=10;move=S}.
             The PowerShell window confirms the server's LAN IP (\texttt{10.50.33.6})
             and the active Wi-Fi adapter (\texttt{central.local}).}
    \label{fig:server_console}
\end{figure}

\begin{figure}[H]
    \centering
    \includegraphics[width=0.92\textwidth]{server_inspector_tcp.png}
    \caption{Unity Inspector showing the \texttt{JPEG Camera Capturer Improved}
             component configured at 1280$\times$720, quality 70. The Console
             confirms a client disconnection event from \texttt{10.50.33.4:41552}
             and the Game view shows the cube displaced by the \texttt{move=S}
             command received from the client.}
    \label{fig:server_inspector}
\end{figure}

\subsection{Client -- Snapshot Window and Terminal}

\begin{figure}[H]
    \centering
    \includegraphics[width=0.85\textwidth]{client_snapshot_tcp.png}
    \caption{Python TCP client (Ubuntu, Lenovo ThinkPad) displaying the Unity
             JPEG snapshot at 1280$\times$720 resolution. The overlay shows:
             frame size 24\,206 bytes, FPS\,=\,10.0, no active movement command
             (\texttt{Move: ???}), and protocol label \texttt{Protocolo: TCP}.
             The terminal confirms the successful connection:
             \texttt{python3 tcp\_client\_viewer.py -ip 10.50.33.6},
             parameters \texttt{w=1280 h=720 q=70 fps=10},
             and \texttt{Conectado a 10.50.33.6:5000}.}
    \label{fig:client_snapshot}
\end{figure}

\begin{figure}[H]
    \centering
    \includegraphics[width=0.85\textwidth]{client_terminal_tcp.png}
    \caption{Python TCP client during an active \texttt{move=S} command. The
             overlay shows frame size 23\,418 bytes, FPS\,=\,10.1, and
             \texttt{Move: S}, causing the cube to move away from the camera
             (visible as a small white square near the horizon). The second
             terminal window demonstrates the \textbf{auto-reconnect} feature:
             after a disconnection (\texttt{Errno\,104 -- Connection reset by
             remote host}), the client automatically reconnects within 2
             seconds (\texttt{Conectado a 10.50.33.6:5000}).}
    \label{fig:client_terminal}
\end{figure}

\subsection{Observed Metrics}

\begin{table}[H]
\centering
\begin{tabular}{@{}lll@{}}
\toprule
\textbf{Parameter} & \textbf{Value} & \textbf{Notes} \\
\midrule
Resolution          & 1280 $\times$ 720 px & HD -- not possible reliably over UDP \\
JPEG quality        & 70               & Configurable via request string \\
Target FPS          & 10               & Server-side pacing \\
Measured FPS        & 10.0             & Stable over Wi-Fi LAN \\
Transport protocol  & TCP              & Persistent connection, no frame loss \\
Server IP           & 10.50.33.6       & Windows\,11, Wi-Fi \texttt{central.local} \\
Client IP           & 10.50.33.4       & Ubuntu, Lenovo ThinkPad \\
Movement commands   & W / A / S / D / 0 & Reset each frame, no drift \\
\bottomrule
\end{tabular}
\caption{System performance metrics during the LAN test.}
\end{table}

% ═══════════════════════════════════════════════════════════════════════════════
\section{Conclusions}

Migrating from UDP to TCP resolves the key limitations encountered in the previous
laboratory session while preserving the same parametric interface and WASD control:

\begin{itemize}
    \item \textbf{No size limit:} 1280$\times$720 JPEG frames are transmitted reliably
          without any IP fragmentation risk.
    \item \textbf{No silent frame loss:} TCP's retransmission mechanism guarantees
          every frame reaches the client, which is critical for a control application.
    \item \textbf{Multi-client support:} each accepted connection spawns an independent
          \texttt{Task}, allowing multiple Python clients to receive the stream
          simultaneously.
    \item \textbf{Auto-reconnect:} the Python client detects disconnection and
          re-establishes the TCP session automatically.
    \item \textbf{Manual framing required:} the 4-byte length prefix adds a small
          implementation overhead not present in UDP, but is a standard pattern for
          TCP message-oriented protocols.
\end{itemize}

\subsection{Future Extensions}

\begin{description}
    \item[Adaptive quality] Dynamically adjust JPEG quality based on measured RTT to
    balance image quality and latency.
    \item[WebSocket transport] Replace raw TCP with WebSockets for browser-based
    clients without a Python dependency.
    \item[H.264 / MPEG streaming] Replace per-frame JPEG with a video codec for
    significantly better compression at high frame rates.
\end{description}

% ── Bibliography ─────────────────────────────────────────────────────────────
\begin{thebibliography}{9}

\bibitem{unity_docs}
Unity Technologies, \textit{Unity Documentation -- Coroutines}, 2024.
\url{https://docs.unity3d.com/Manual/Coroutines.html}

\bibitem{msdn_tcp}
Microsoft, \textit{TcpListener Class -- .NET API Reference}, 2024.
\url{https://learn.microsoft.com/en-us/dotnet/api/system.net.sockets.tcplistener}

\bibitem{msdn_networkstream}
Microsoft, \textit{NetworkStream Class -- .NET API Reference}, 2024.
\url{https://learn.microsoft.com/en-us/dotnet/api/system.net.sockets.networkstream}

\bibitem{opencv_python}
OpenCV Team, \textit{OpenCV-Python Tutorials}, 2024.
\url{https://docs.opencv.org/4.x/d6/d00/tutorial_py_root.html}

\bibitem{tutorial}
R. Marín Prades, \textit{Tutorial JPEG Capture in Unity and UDP Communication
for Remote Control}, February 2026.

\end{thebibliography}

\end{document}
